% Answer key for Chapters 16--23. This file is \input directly inside the document.
% Every number comes from the Python test scripts (see the report).

\section*{Chapter 16: Differential Equations and Dynamical Systems}
\begin{enumerate}[label=\textbf{16.\arabic*}]
  \item $y=Ce^{t^{2}}$ with $y(0)=1$ gives $y=e^{t^{2}}$, so $y(1)=e=2.718282$.
        Euler with $h=0.1$ uses $y_{n+1}=(1+0.02n)y_{n}$; the result is $2.334633$,
        error $3.8365\times10^{-1}$. That method lags behind because it uses
        the slope at the left end of each step.

  \item The difference $w=y_{1}-y_{2}$ satisfies $w'=-p(t)w$ and $w(t_{0})=0$. Multiply
        by $e^{\int p}$: $\bigl(we^{\int p}\bigr)'=0$, so $we^{\int p}$ is constant, and
        its value is zero at $t_{0}$. Hence $w\equiv0$. No solution formula is used,
        only the product rule and uniqueness of initial values.

  \item Not enough evidence. Linearization sees only the degree-one terms; the
        eigenvalues $\pm i$ say only that the linear part is a rotation. With $r'=r^{3}$
        the origin is a center for the linearization but an outward spiral, whereas
        $r'=-r^{3}$ gives it an inward spiral. So $\pm i$ does not determine stability.

  \item $E'=vv'+\omega_{0}^{2}\theta\theta'=v(-\omega_{0}^{2}\theta)+\omega_{0}^{2}\theta
        v=0$, so $E$ is constant. The equation $v^{2}/(2E)+\theta^{2}/(2E/\omega_{0}^{2})=1$
        is an ellipse, and $\theta=A\cos(\omega_{0}t+\varphi)$ closes one turn
        in $2\pi/\omega_{0}$.

  \item $\tau=RC=2.2\times10^{3}\cdot10^{-5}=0.022$ seconds. At $t=\tau$:
        $V=12(1-e^{-1})=7.585447$ volts. For $6$ volts, $t=-\tau\ln(1-\tfrac{6}{12})
        =0.022\cdot0.693147=15.249$ milliseconds.

  \item With $v=2\cos(\theta/2)$: $v^{2}/2=2\cos^{2}(\theta/2)$ and
        $1-\cos\theta=2\sin^{2}(\theta/2)$, and their sum is $E=2$. At $\theta=\pm\pi$
        the value is $v=0$, so the curve joins those two points; it separates
        oscillation from full rotation, and the time needed to travel it is infinite.

  \item $T_{0}=2\pi\sqrt{0.25/9.80665}=1.003205$ seconds.
        $T=T_{0}\cdot\tfrac{2}{\pi}K(\sin15^{\circ})=1.020669$ seconds; the difference is
        $1.74$ percent. An amplitude of $30^{\circ}$ lengthens the period, it does not
        shorten it.

  \item The factor per step is $1+h\lambda=1-100h$; the condition $|1-100h|<1$ gives
        $0<h<0.02$. Exactly at $h=0.02$ the factor is $-1$: the value changes sign
        with fixed magnitude $1$, neither decaying nor blowing up. That is the
        marginal stability boundary, not stability.

  \item Suppose $y$ is periodic and not constant. At its maximum point $y'=0$, so
        $f(y_{\max})=0$. The constant solution $y\equiv y_{\max}$ passes through that
        point, and uniqueness (Lipschitz) forces the solutions to be the same,
        contradicting non-constancy. So there is no non-constant periodic orbit.
\end{enumerate}

\section*{Chapter 17: Probability and Conditional Probability}
\begin{enumerate}[label=\textbf{17.\arabic*}]
  \item 
$P(A\cup B\cup C)=P(A\cup B)+P(C)-P\bigl((A\cup B)\cap C\bigr)$, then
        $P\bigl((A\cup B)\cap C\bigr)=P(A\cap C)+P(B\cap C)-P(A\cap B\cap C)$. The
        triple-overlap term is subtracted three times, once for each pair, and then
        added back once with a $+$ sign.

  \item False on general spaces. $P(A\mid B)=0$ only means $P(A\cap B)=0$. On a
        finite space with positive probabilities, that forces $A\cap B=\varnothing$.
        On a continuous space it does not: $A=\{1/2\}$ and $B=(0,1)$ uniform give
        $P(A\cap B)=0$ even though $1/2\in A\cap B$.

  \item From $10\,000$ people: $200$ are sick, $190$ correctly positive; with specificity
        $0.90$ there are $980$ false positives, so $190/1170=19/117\approx0.1624$.
        With $0.99$: only $98$ false positives remain, so $190/288=95/144
        \approx0.6597$: the precision rises fourfold.

  \item $P(A)=P(A\cap B)+P(A\cap B^{c})$, so
        $P(A\cap B^{c})=P(A)-P(A)P(B)=P(A)(1-P(B))=P(A)P(B^{c})$. For
        $A^{c},B^{c}$ use $P(A^{c}\cap B^{c})=1-P(A\cup B)$ and inclusion--exclusion,
        or apply the result above twice.

  \item $1-\prod_{k=0}^{29}(1-k/365)=0.706316$, whereas
        $1-e^{-435/365}=0.696320$; the difference is $0.009996$. The approximation uses
        $1-x\approx e^{-x}$ for $x$ up to $29/365$, and treats pair events
        as mutually independent when they are not.

  \item $A=(A\cap B)\cup(A\cap B^{c})$ and the two pieces are disjoint, because
        $B\cap B^{c}=\varnothing$. Additivity gives
        $P(A)=P(A\cap B)+P(A\cap B^{c})=P(A\mid B)P(B)+P(A\mid B^{c})P(B^{c})$,
        using the definition of conditional probability twice.

  \item The probability of winning by switching is $3/4=0.75$: the initial choice wins
        only when the car really is there ($1/4$). A simulation of $10^{4}$ games gives
        about $0.75$ with standard deviation $\sqrt{0.75\cdot0.25/10^{4}}\approx0.0043$.

  \item 
The correct figure is $0{,}9075$, not $0{,}9999$: with $0{,}1$ percent
        prevalence only $1$ person in $1000$ is ill, while the false positives among
        two positive tests are $999\cdot0{,}01^{2}\approx0{,}1$. Hence
        $0{,}99^{2}\cdot0{,}001/(0{,}99^{2}\cdot0{,}001+0{,}01^{2}\cdot0{,}999)=0{,}9075$.
        The assumptions: the two tests are conditionally independent and both
        sensitivity and specificity are exactly $0{,}99$.

  \item Not enough. Take $A$ the first coin shows heads, $B$ the second coin shows
        heads, $C$ both coins match. Every pair gives $1/4=P\cdot P$, but
        $P(A\cap B\cap C)=1/4$ whereas $P(A)P(B)P(C)=1/8$; the difference is $1/8$.
\end{enumerate}

\section*{Chapter 18: Random Variables and Distributions}
\begin{enumerate}[label=\textbf{18.\arabic*}]
  \item $c=1-(0.1+0.4)=0.5$; $\mathbb{E}[X]=-0.2+0+1.5=1.3$;
        $\mathbb{E}[X^{2}]=0.4+0+4.5=4.9$; $\operatorname{Var}(X)=4.9-1.69=3.21$.

  \item $(-\infty,b]=(-\infty,a]\cup(a,b]$ and the two pieces are disjoint, so
        $F(b)=F(a)+P(a<X\le b)$. For discrete variables $P(X=a)>0$ makes
        $P(a\le X\le b)$ larger; for continuous ones $P(X=a)=0$ so the two agree.

  \item It is true that $P(X=x)=0$ for every $x$ in a continuous distribution. The
        mistake is concluding ``impossible'': probability zero does not mean the empty
        set. Countable additivity gives probability zero to a countable union of points,
        but an interval contains uncountably many points and has probability one.

  \item Differentiate $\sum q^{k}=1/(1-q)$ twice: $\sum k(k-1)q^{k-2}=2/(1-q)^{3}$.
        Since $\sum k^{2}q^{k-1}=(1+q)/(1-q)^{3}$, with $q=1-p$ we get
        $\mathbb{E}[X^{2}]=p(1+q)/p^{3}=(1+q)/p^{2}=(2-p)/p^{2}$.
        Hence $\operatorname{Var}=(2-p)/p^{2}-1/p^{2}=(1-p)/p^{2}$; for $p=0.2$:
        $\mathbb{E}[X^{2}]=45$ and $\operatorname{Var}=20$.

  \item $\sigma=2$: $P(X>1)=1-\Phi(-1)=0.841345$;
        $P(1<X<5)=2\Phi(1)-1=0.682689$; the $0.2$ quantile is
        $3+2(-0.841621)=1.316758$. All three agree with
        \texttt{scipy.stats.norm} to six digits.

  \item $\sum_{k=0}^{40}e^{-2.5}2.5^{k}/k!=1-5.39\times10^{-35}$, that is $1$
        to every printed digit. $P(X\ge8)=0.004247$, whereas the binomial with
        $n=1000$, $p=0.0025$ gives $0.004191$; the binomial is smaller because
        its variance is $2.4938<2.5$ and it is bounded above by $n$.

  \item $P(X>s+t\mid X>s)=e^{-\lambda(s+t)}/e^{-\lambda s}=e^{-\lambda t}$.
        Geometric: $P(X>k)=q^{k}$, so the ratio is $q^{s+t}/q^{s}=q^{t}$. Uniform on
        the integers fails because the tail probability runs out, and the ratio depends
        on $s$.

  \item $f_{Y}(y)=\frac{1}{4\sqrt{y}}$ for $0<y<1$ (two preimages) and
        $\frac{1}{8\sqrt{y}}$ for $1\le y<9$ (one preimage). The integral is
        $\tfrac12+\tfrac12=1$. $\mathbb{E}[Y]=\frac14\int_{-1}^{3}x^{2}\,dx=7/3
        =2.3333$, the same as $\mathbb{E}[X^{2}]$, and via $f_{Y}$:
        $\int_{0}^{1}\frac{y}{4\sqrt{y}}\,dy+\int_{1}^{9}\frac{y}{8\sqrt{y}}\,dy
        =\tfrac16+\tfrac{13}6=\tfrac73$.

  \item Take $X$ uniform on $\{-1,0,1\}$ and $Y=|X|$. Then
        $\mathbb{E}[XY]=\mathbb{E}[X]=0$, so $\operatorname{Cov}=0$ and $\rho=0$.
        But $P(X=0,Y=0)=1/3$ whereas $P(X=0)P(Y=0)=1/9$: they are not independent.

  \item 
$\operatorname{se}(\bar X)=\sqrt{p(1-p)/(nR)}\le0{,}005$ gives
        $R\ge0{,}2275/(20\cdot0{,}000025)=455$. For the variance,
        $\operatorname{se}(\bar s^{2})=\sigma^{2}\sqrt{2/((n-1)R)}$ (normal
        approximation) with $\sigma^{2}=4{,}55$ and $n=20$, so
        $R\ge2\cdot4{,}55^{2}/(19\cdot0{,}02^{2})=5\,448{,}03$, that is $R=5\,449$.
        Using $2\cdot4{,}55^{2}/0{,}02^{2}$ alone drops the factor $n-1$ and gives a
        figure nineteen times too large.
  \end{enumerate}

\section*{Chapter 19: The Law of Large Numbers and the Central Limit Theorem}
\begin{enumerate}[label=\textbf{19.\arabic*}]
  \item $\{|(X_{n}+Y_{n})-(X+Y)|>\varepsilon\}$ is contained in
        $\{|X_{n}-X|>\varepsilon/2\}\cup\{|Y_{n}-Y|>\varepsilon/2\}$. The union bound
        gives a probability of at most the sum of two terms, each going to $0$.
        The $\varepsilon/2$ is needed so that two errors both below it add up to
        less than $\varepsilon$.

  \item $\mathbb{E}[X]=0$ and $\mathbb{E}[X^{2}]=2a^{2}\cdot\sigma^{2}/(2a^{2})=\sigma^{2}$,
        so $\operatorname{Var}(X)=\sigma^{2}$. Then
        $P(|X|\ge a)=\sigma^{2}/a^{2}$, exactly equal to the Chebyshev bound: the bound
        is attained, so it is sharp.

  \item False. The central limit theorem demands finite variance and iid observations; it
        also speaks only about the limit, not about $n=30$. In the lab section, the
        exponential with $n=30$ still gives $D=0.0219$ with $p=0.000$, and Cauchy
        gives a median $|\bar X_{n}|$ that settles near $1$.

  \item Chebyshev: $n\ge\sigma^{2}/(\varepsilon^{2}\delta)=0.25/(0.0025\cdot0.05)
        =2000$. The exact binomial computation gives $n=371$ (tail $0.0484$). The
        Chebyshev bound is five times looser.

  \item $\mathbb{E}[X_{n}]=n^{2}\cdot(1/n)=n\to\infty$, whereas for fixed
        $\varepsilon$ we have $P(|X_{n}|>\varepsilon)=1/n\to0$. So $X_{n}\to0$ in
        probability but its expectation explodes: no domination is available.

  \item Exact: $0.069545$. Without the correction, $z=(50-60)/\sqrt{42}=-1.5430$ gives
        $0.061411$, error $0.008134$. With the correction at $50.5$:
        $0.071340$, error $0.001795$. The continuity correction improves it fourfold.

  \item Median $|\bar X_{n}|$: $1.004$; $1.008$; $1.040$; $0.973$ for
        $n=10$ up to $10^{4}$. $P(|\bar X_{n}|>10)$: $0.0630$; $0.0680$;
        $0.0653$; $0.0580$, theory $0.0635$. The standard deviation $33.0$; $863.6$;
        $23.9$; $379.5$ does not shrink because the Cauchy variance is infinite.

  \item $\varphi_{\bar X_{n}}(t)=\bigl[\varphi(t/n)\bigr]^{n}
        =\bigl[e^{-|t|/n}\bigr]^{n}=e^{-|t|}$, so the average is again Cauchy for
        every $n$. Hence the se stays as large as the scale of the data itself, and
        ``mean $\pm$ two se'' does not narrow even as $n$ grows.

  \item $n=(1.96\cdot0.491971/0.0005)^{2}=3\,719\,214$; for half-width
        $0.001$ it is enough to take $929\,803$. The ratio is $4$, because $n\propto h^{-2}$:
        halve the width and the sample size rises fourfold.

  \item For $t<c$ and $\varepsilon=c-t$:
        $P(X_{n}\le t)\le P(X_{n}\le c-\varepsilon)+P(|X_{n}-c|>\varepsilon)\to0$;
        for $t>c$ use the complement. The constant condition cannot be dropped:
        take $X$ Rademacher and $X_{n}=(-1)^{n}X$. Every $X_{n}$ is still
        Rademacher (because $X$ is symmetric), so $X_{n}\rightsquigarrow X$; but for
        odd $n$ we have $|X_{n}-X|=2$ with probability $1$, so $X_{n}$ does not
        converge to $X$ in probability.
\end{enumerate}

\section*{Chapter 20: Estimation, Confidence Intervals, and Hypothesis Testing}
\begin{enumerate}[label=\textbf{20.\arabic*}]
  \item False. A p-value of $0.42$ means the data are not enough to contradict $p=0.5$;
        it does not prove $p=0.5$. The supplementary computations: a $95$ percent
        interval, for instance $(0.28;\,0.57)$, and the power of the test at the
        difference considered to matter.

  \item $\mathbb{E}[\hat p(1-\hat p)]=p-\bigl(p^{2}+p(1-p)/n\bigr)
        =p(1-p)(n-1)/n$. The unbiased estimator is
        $\frac{n}{n-1}\hat p(1-\hat p)$. For $p=0.3$, $n=20$: the expectation is
        $0.21\cdot19/20=0.1995$, and that estimator brings it back to $0.21$.

  \item Write $\hat\theta_{n}-\theta=\operatorname{bias}_{n}
        +(\hat\theta_{n}-\mathbb{E}\hat\theta_{n})$. The first term goes to $0$; the
        second goes to $0$ in probability by Chebyshev, because its variance goes to
        $0$. The MSE $=\operatorname{bias}^{2}+\operatorname{Var}$ going to $0$ is stronger,
        but only the two conditions above are used.

  \item 
$7/40$: Wilson $(0{,}0875;\,0{,}3195)$, width $0{,}2320$; Wald
        $(0{,}0572;\,0{,}2928)$, width $0{,}2355$. $13/20$: Wilson width $0{,}3860$,
        Wald $0.4181$. The $7/40$ interval is narrower, but that is not a sign of
        higher precision. Partly it is $n=40$; partly it is $\hat p=0.175$ being far
        from $0.5$, which shrinks $\hat p(1-\hat p)$: at $n=40$ with $\hat p=0.5$ the
        Wilson width is $0.2960$, still wider than $0.2320$. (At equal $n$, $7/20$ and
        $13/20$ both give $0.3860$, since $\hat p(1-\hat p)$ is symmetric about
        $1/2$.) What settles precision is not width but coverage: the Wald coverage
        at $p=0.175$, $n=40$ is only $0.921$, whereas Wilson gives $0.943$. A narrow
        interval near the edge is paid for with coverage below $95$ percent.

  \item $|k-15|>1.96\sqrt{7.5}=5.3676$, so reject when $k\le9$ or $k\ge21$.
        Its exact size is $0.042774$; its power at $p=0.7$ is only $0.588816$.
        With $n=100$ (reject when $k\le40$ or $k\ge60$) the power rises to $0.987502$.

  \item Wald: $|\hat p-p_{0}|\le z\sqrt{\hat p(1-\hat p)/n}$, which is equivalent to
        $p_{0}$ lying in $\hat p\pm z\sqrt{\hat p(1-\hat p)/n}$. A test that uses
        $\sqrt{p_{0}(1-p_{0})/n}$ uses a different denominator, so the set of $p_{0}$
        that are not rejected is not that Wald interval.

  \item 
Suppose the observed difference is $0.87$ hours with standard error $0.86$ hours.
``At the $5$ percent level the data do not reject equality of the two methods
        ($p=0.31$). The $95$ percent interval for the difference is
        $(-0.82;\,2.56)$ hours, so a $2$-hour difference is still a plausible value
        and even a $0.5$-hour difference cannot be ruled out. The power of this test
        against a difference of $0.5$ hours is only $0.09$: absence of evidence is
        not evidence of absence, and the honest conclusion is that these data do not
        decide the question.''

  \item Rule of three: $3/100=0.03$. The exact one-sided bound is
        $1-0.05^{1/100}=0.029513$; the Wilson upper bound is $0.036993$. The three
        differ because the rule of three is only an approximation and Wilson is two-sided.
        For the statement ``at most $5$ percent'' to be valid at the $95$ percent level:
        the exact bound is $0.0295$.

  \item $\bar x=19.125$, $s=1.082191$, $\operatorname{se}=0.312401$;
        $t_{11;0.95}=1.795885$ gives $(18.5640;\,19.6860)$ with width
        $1.1221$, whereas $95$ percent gives width $1.3752$. The width scales as
        $1/\sqrt{n}$, so halving it needs $n=48$.
\end{enumerate}

\section*{Chapter 21: Regression and Applied Statistics}
\begin{enumerate}[label=\textbf{21.\arabic*}]
  \item Not enough evidence. $R^{2}=0.99$ says only that the variance is explained;
        $s$ determines the error. On the cafe data $R^{2}=0.992821$ but
        $s=0.365148$ hours $=21.9$ minutes, and the half-width of the $95$ percent
        prediction interval at $x_{0}=5$ is $0.9101$ hours, that is $55$ minutes. What
        should be quoted: $s$ and a prediction interval.

  \item The normal equations are $X^{\top}e=0$. The first column of $X$ holds ones, so
        $\sum e_{i}=0$; the second column holds $x_{i}$, so $\sum x_{i}e_{i}=0$.
        Without the ones column (a line through the origin), $\sum e_{i}=0$ is no longer
        imposed and generally fails.

  \item $\hat b-\beta=\sum_{i}(x_{i}-\bar x)\varepsilon_{i}/S_{xx}$. The terms are
        independent and $\operatorname{Var}(\varepsilon_{i})=\sigma^{2}$, so
        $\operatorname{Var}(\hat b)=\sigma^{2}\sum(x_{i}-\bar x)^{2}/S_{xx}^{2}
        =\sigma^{2}/S_{xx}$. No matrix formula is used.

  \item $\hat b=0.7$, $\hat a=1.3$, $S_{xx}=10$, $s=0.516398$,
        $\operatorname{se}(\hat b)=0.163299$; with $t_{3}=3.182446$ the interval is
        $(0.1803;\,1.2197)$. It is wider than for the cafe data because $n=5$,
        $S_{xx}$ is small, and $s$ is large.

  \item 
In the cafe data $\bar x=5$, $n=9$, $S_{xx}=60$, so the leverage is
        $c(x_{0})=1/9+(x_{0}-5)^{2}/60$: $c(2)=0.261111$ and $c(9)=0.377778$.
        With $2.364624\sqrt{0.133333\,(1+c)}$ the half-widths are $0.9696$ for
        $x_{0}=2$ and $1.0135$ for $x_{0}=9$; they differ because
        $(x_{0}-\bar x)^{2}/S_{xx}$ is larger farther from the centre. That term
        shrinks like $1/n$, so the difference between the two intervals does vanish
        for large $n$: both tend to $1.96\sigma=0.716$ half-width. What never
        vanishes is the width itself, because the $+1$ term (the error of one new
        observation) remains: a prediction interval does not shrink to zero, it
        shrinks to the spread of the error.

  \item The column space of the big model contains the column space of the small model, so
        its projection is closest: the SSE cannot rise. $R^{2}$ may stay the same, but the
        divisor $n-p$ shrinks, so $R^{2}_{\text{adj}}$ falls: from $0.991795$ ($p=2$) to
        $0.990428$ ($p=3$).

  \item $\mathrm{VIF}=1/(1-0.81)=5.263158$, so
        $\operatorname{se}(\hat b_{1})$ grows by a factor $\sqrt{5.263158}$. With
        $x_{2}$: $\sqrt{2/(50\cdot0.19)}=0.458831$. Without $x_{2}$:
        $\sqrt{2/50}=0.2$. The ratio is $2.294157$.

  \item 
Group A: $(1;2)$, $(2;1{,}8)$, $(3;1{,}6)$; group B: $(5;4{,}0)$,
        $(6;3{,}8)$, $(7;3{,}6)$. Each group has slope $-0{,}2$, the pooled slope is
        $+0{,}4$. The valid decomposition adds contributions, not slopes:
        $b_{\text{gab}}=(S_{xy}^{\text{dalam}}+S_{xy}^{\text{antar}})/S_{xx}=-0{,}8/28+12/28=-0{,}028571+0{,}428571=0{,}4$,
        whereas the between-group slope on its own is $12/24=0{,}5$.

  \item ``Of the twelve models tried, the best gives $p=0.03$. After
        correcting for twelve comparisons (Bonferroni: $0.36$) the effect is no longer
        significant; therefore these data have not decided.''
\end{enumerate}

\section*{Chapter 22: Floating-Point Arithmetic and Error Analysis}
\begin{enumerate}[label=\textbf{22.\arabic*}]
  \item $\operatorname{ulp}(1)=2^{-52}=2.220\times10^{-16}$;
        $\operatorname{ulp}(1024)=2^{-42}=2.274\times10^{-13}$;
        $\operatorname{ulp}(0.1)=2^{-56}=1.388\times10^{-17}$;
        $\operatorname{ulp}(3\times10^{8})=2^{-24}=5.960\times10^{-8}$.
        $|x|u$ differs from the ulp by at most a factor of two, because the ulp is constant
        inside one binade. In $[1,2)$ there are $2^{52}$ numbers.

  \item $2^{-53}$ is exactly half an ulp of $1$, and rounding to even sends it back
        to $1$: $1+2^{-53}=1$. Since $2^{-52}$ is one full ulp,
        $1+2^{-52}>1$. In binary32, $u=2^{-24}$ and ulp$(1)=2^{-23}$, so the smallest
        positive number is $2^{-23}=1.192\times10^{-7}$.

  \item $\operatorname{fl}(x)$ is the nearest number, so
        $|\operatorname{fl}(x)-x|\le\tfrac12\operatorname{ulp}(x)$ by the definition of
        rounding. Equality cannot hold for all $x$: when $x$ is itself representable the
        error is $0$, and $\tfrac12$ ulp is attained only halfway between two numbers.

  \item False. Every operation is rounded correctly, but the errors accumulate along the
        algorithm. Example from this chapter: the sequential sum $1+10^{6}\cdot10^{-16}$
        stays $1$, whereas Kahan gives $1.0000000001$; the difference $10^{-10}$
        is far above $u$.

  \item With denominator $2^{54}$: $(0.1+0.2)+0.3=10808639105689192/2^{54}$ and
        $0.1+(0.2+0.3)=10808639105689190/2^{54}$. The difference is $2/2^{54}=2^{-53}$,
        exactly one ulp. \texttt{fsum} gives $0.6$; \texttt{float.hex} differs in the
        last digit.

  \item The naive formula gives $0.4996003611$; the double-angle
        identity gives $0.4999999999999996$ and $0.5-x^{2}/24+x^{4}/720$
        agrees. Its relative error is $-7.99\times10^{-4}$ and comes from
        cancellation in $1-\cos x$: $\cos x$ discards the digits the difference needs.

  \item Naive: the small root is $-1.0000000111\times10^{-4}$. The stable form
        $-2c/\bigl(b+\sqrt{b^{2}-4ac}\bigr)$ gives
        $-1.0000000100\times10^{-4}$, relative error $1.12\times10^{-9}$.
        The condition number of the subtraction is $\bigl(2\cdot10^{4}\bigr)/\bigl(2\cdot
        10^{-4}\bigr)=10^{8}$, so about eight digits are lost.

  \item At level $j$ there are $2^{k-j}$ additions, each with error at most
        $u$ times the sum of the absolute values of its inputs; summed over the whole tree,
        the error is $\le ku\sum|x_{i}|$ with $k=\log_{2}n$. That tree depth is what
        replaces $n$ in the bound for the sequential sum.

  \item Sequential: every $10^{-16}$ is below $\tfrac12\operatorname{ulp}(1)
        =1.11\times10^{-16}$, so the result stays $1$. Kahan:
        $1.0000000001$, the same as the exact value. Each Kahan step is also rounded,
        but the lost remainder is kept in $c$ and returned on the next step.

  \item $\kappa\approx1+\bar x^{2}/s^{2}$. For $1000$ values
        $999+\bigl((k\bmod7)-3\bigr)/4$ with $k=1,\dots,1000$:
        $\bar x=999.00075$, $s=0.4997$ (population standard deviation, matching
        the $1/n$ divisor in the problem),
        $\kappa\approx4.0\times10^{6}$. With $u_{32}=5.96\times10^{-8}$,
        $\kappa u\approx0.238$: almost all digits are lost. Measured: one pass
        errors by $25$ percent, two passes by $4.4\times10^{-8}$.
\end{enumerate}

\section*{Chapter 23: Convex Optimization and Duality}
\begin{enumerate}[label=\textbf{23.\arabic*}]
  \item $f$: Hessian $\begin{pmatrix}2&1\\1&2\end{pmatrix}$, eigenvalues $1$ and $3$,
        positive definite, strictly convex. $g$: $\begin{pmatrix}0&1\\1&0\end{pmatrix}$,
        eigenvalues $\pm1$, indefinite, neither of the two. $h$:
        $\begin{pmatrix}4&-4\\-4&6\end{pmatrix}$, eigenvalues $0.8769$ and $9.1231$,
        positive definite. None is singular semidefinite; for $x_{1}^{2}+2x_{1}x_{2}
        +x_{2}^{2}$ (eigenvalues $0$ and $4$) the minimum is not unique.

  \item For $x,y\in C_{1}\cap C_{2}$ and $t\in[0,1]$, the mixed point lies in
        $C_{1}$ and in $C_{2}$, hence in their intersection. Union: $[0,1]\cup[2,3]$ contains
        $1$ and $2$ but not $(1+2)/2=1.5$.

  \item $f_{i}(tx+(1-t)y)\le tf_{i}(x)+(1-t)f_{i}(y)\le tM(x)+(1-t)M(y)$ for each
        $i$, then take the maximum. For $\max\{x_{1}^{2},x_{2}\}$ at $(0,0)$ the gradient
        does not exist (a crease along $x_{2}=x_{1}^{2}$), but its subdifferential
        $\operatorname{conv}\{(0,0),(0,1)\}$ contains $\mathbf{0}$: a minimum. At $(1,1)$
        the subdifferential is $\operatorname{conv}\{(2,0),(0,1)\}$ with no zero, not a
        minimum.

  \item Without $x_{1}\ge1$: $x_{1}=0.8$, $x_{2}=1.6$, which violates it. So $x_{1}=1$,
        $x_{2}=1.5$, $p^{*}=3.25$. The multipliers are $\nu=-1.5$ for the equality and
        $\lambda=0.5\ge0$ for the inequality. \texttt{scipy.optimize.mini\-mize}
        gives the same value.

  \item $g(\lambda)=\min_{x}\bigl[x_{1}^{2}+x_{2}^{2}+\lambda(2-x_{1}-x_{2})\bigr]
        =2\lambda-\lambda^{2}/2$. Its derivative is $2-\lambda$, so the maximum is at
        $\lambda=2$ with $d^{*}=2=p^{*}$. The gap is zero because the constraint is affine
        and $x=(1.5;\,1.5)$ is strictly feasible: Slater holds.

  \item False. Strong duality only equates the values, it does not guarantee attainment.
        Example: minimize $e^{x}$ over $\mathbb{R}$. $p^{*}=0$ but it is not attained;
        the dual side is also $0$, so the gap is zero while the primal optimum does not
        exist.

  \item False. Complementary slackness states $\lambda_{i}g_{i}(\mathbf{x}^{*})=0$:
        an active constraint may have a zero multiplier, and an inactive constraint must
        have a zero multiplier. The implication runs one way, so ``active $\Rightarrow$
        positive'' is not valid.

  \item Without the bound, $(x_{1}-2)^{2}+(x_{2}-2)^{2}$ on $x_{1}+x_{2}=2$ is minimized
        at $x_{1}=x_{2}=1$. That satisfies $x_{1}\le c$ if and only if $c\ge1$, so
        at $c=1$ the multiplier is $\lambda=0$ and $p^{*}=2$; the constraint starts to
        bind only for $c<1$ (at $c=0.5$: $p^{*}=2.5$, $\lambda=2$).

  \item For $\lambda>0$ and $x_{k}\to-\infty$, the term $\lambda x_{k}^{3}$ dominates
        and is negative, so $g(\lambda)=-\infty$. Only $\lambda=0$ is finite:
        $g(0)=\min(x-2)^{2}=0$, hence $d^{*}=0$. Meanwhile $p^{*}=1$ at $x=1$ (the piece
        $[0,1]$), because at $x=-1$ the value is $9$. The gap is $1$, one hundred percent.

  \item The first-order inequality for a convex function gives
        $f(y)\ge f(x^{*})+\nabla f(x^{*})^{\top}(y-x^{*})=f(x^{*})$ for every $y$,
        without using strict convexity. That proof fails for $f(x)=x^{3}$ because
        $f$ is not convex; the condition violated is precisely convexity of $f$, at $x=0$
        where $\nabla f=0$ but it is not a minimum.
\end{enumerate}
